\hypertarget{appendix-942ad41afb30}{}\label{referencestart:appendix-942ad41afb30}
\section{Compiler IR and machine instructions}\label{appendix-942ad41afb30-compiler-ir-and-machine-instructions}

\hypertarget{appendix-942ad41afb30-compiler-ir-and-machine-instructions}{}

\textbf{Curriculum depth:} Forward · \textbf{First seen:} C001 · \textbf{Latest development:} C002

\begin{allmightroledefinition}

\subsection{Definition and mechanism}\label{appendix-942ad41afb30-definition-and-mechanism}

\hypertarget{appendix-942ad41afb30-definition-and-mechanism}{}

Source-level operators are compiled through intermediate representation into target-specific instructions and physical state changes.

\end{allmightroledefinition}

IR is the compiler's intermediate representation of program operations. Instruction selection translates it into target instructions under the source-language contract. A hardware add may retain low bits, while an optimizer can assume a signed C++ addition never overflows.

\subsection{Prerequisites}\label{appendix-942ad41afb30-prerequisites}

\hypertarget{appendix-942ad41afb30-prerequisites}{}

\begin{itemize}
\tightlist
\item
  \hyperref[appendix-f414bd4ed05c]{Interpretation contract} --- For \BookInlineCode{{0}{x}{A}{D}}{0xAD}, unsigned interpretation gives 173; eight-bit two's complement gives −83.
\item
  \hyperref[appendix-d67a19507dc5]{Machine word} --- A processor's natural integer processing width is often called its word width, but the term is architecture dependent.
\end{itemize}

\subsection{Used by}\label{appendix-942ad41afb30-used-by}

\hypertarget{appendix-942ad41afb30-used-by}{}

\begin{itemize}
\tightlist
\item
  \href{../01-Bits-Programs-and-Reliable-Systems-P02/main.pdf\#nameddest=appendix-806686359e1b}{Translation units and linking (Bits, Programs, and Reliable Systems, Part 2, Reference notes: Translation units and linking)} --- uses this prerequisite to establish its entity, lifetime, representation or lookup contract.
\end{itemize}

Connect later chunks here when they use this concept.

\subsection{Recall and next depth}\label{appendix-942ad41afb30-recall-and-next-depth}

\hypertarget{appendix-942ad41afb30-recall-and-next-depth}{}

Explain compiler ir and machine instructions without the label, work the example above, then state the assumption whose removal would change the result.

\textbf{Pending:} lowering, instruction selection, optimization, and microarchitecture.

\subsection{Source and navigation}\label{appendix-942ad41afb30-source-and-navigation}

\hypertarget{appendix-942ad41afb30-source-and-navigation}{}

\hyperref[01-foundations-i-ch001-from-a-signal-to-an-interpretation]{C001 teaching chapter} · \hyperref[01-foundations-i-ch001-arithmetic-becomes-a-policy]{C002 development} · \href{../01-Bits-Programs-and-Reliable-Systems-P03/main.pdf\#nameddest=appendix-fd090d982a15}{Part definitions and mechanisms (Bits, Programs, and Reliable Systems, Part 3, Reference notes: Foundations)} · \href{../07-Vision-XR-and-Playable-Worlds-P01/main.pdf\#nameddest=appendix-4ad47bdd269a}{Dependency Map (Vision, XR, and Playable Worlds, Part 1, Reference notes: Dependency Map)}
